\chapter{Comparaci\'on T\'ecnica}

A continuaci\'on se presenta una comparaci\'on entre distinto sistemas operativos, para ello considerando la arquitectura del kernel, lenguaje de implementaci\'on, soporte de hardware, documentaci\'on y nivel de complejidad.

\begin{landscape}
% Poner esta tabla al final como resumen - nomas este es un ejemplo de uso 
\begin{muntab}{lll}{comparacionSO1}{Comparación técnica de sistemas operativos - Arquitectura y lenguaje}
\hline
\textbf{Sistema} & \textbf{Arquitectura} & \textbf{Lenguaje} \\
\hline
GNU/Linux & Monolítico modular (módulos cargables) & C, ASM; Rust (drivers recientes) \\
FreeBSD & Monolítico con modularidad & C, ASM \\
Minix 3 & Microkernel (servicios en userland) & C, ASM \\
Redox OS & Microkernel (drivers en userland) & Rust \\
Microsoft Windows & Híbrido (NT: Executive + Kernel) & C, C++, ASM; C\# en capas altas \\
SerenityOS & Monolítico estilo Unix & C++20 (lib propia) \\
FreeRTOS & Kernel RTOS monolítico, preemptivo & C (port específico con ASM) \\
macOS (Darwin/XNU) & Híbrido (Mach + BSD + I/O Kit) & C (Mach/BSD), C++ (I/O Kit) \\
Haiku OS & Híbrido (derivado de NewOS) & C++ (API/GUI); C \\
Android (AOSP) & Kernel Linux monolítico + framework & C/C++ (kernel/nativo), Java/Kotlin (framework) \\
\hline
\end{muntab}
\end{landscape}

\begin{landscape}
\begin{muntab}{llll}{comparacionSO2}{Comparación técnica de sistemas operativos - Tamaño, Kernel y soporte}
\hline
\textbf{Sistema} & \textbf{Tamaño} & \textbf{Kernel} & \textbf{Soporte} \\
\hline
GNU/Linux & Muy grande & Linux & Comunitario + empresas; LTS \\
FreeBSD & Grande & FreeBSD & Core Team / FreeBSD Foundation \\

Minix 3 & Pequeño/Mediano & MINIX (microkernel) & Académico/comunitario \\

Redox OS & Mediano & Redox (microkernel) & Comunitario (GitHub) \\

Microsoft Windows & Muy grande & Windows NT & Comercial (Microsoft) \\

SerenityOS & Mediano & Serenity & Comunitario (GitHub) \\

FreeRTOS & Muy pequeño & FreeRTOS & MIT; soporte AWS (LTS) \\

macOS (Darwin/XNU) & Muy grande & XNU & Comercial (Apple); base Darwin abierta \\

Haiku OS & Mediano & Haiku & Comunitario (Haiku, Inc.) \\

Android (AOSP) & Muy grande & Linux (ramas Android) & Google/OEMs (AOSP + CDD) \\
\hline
\end{muntab}
\end{landscape}

\begin{landscape}
\begin{muntab}{lll}{comparacionSO3a}{Comparación técnica de sistemas operativos - HW (Arquitecturas) y documentación}
\hline
\textbf{Sistema} & \textbf{HW (arquitecturas)} & \textbf{Documentación} \\
\hline
GNU/Linux & x86/x86\_64, ARM/ARM64, RISC-V, PowerPC, más & kernel.org; TLDP; guías de distros \\

FreeBSD & x86/x86\_64, ARM/ARM64, PowerPC, RISC-V & FreeBSD Handbook; man pages \\

Minix 3 & x86/x86\_64 (ports limitados) & minix3.org; Design Doc; man \\

Redox OS & x86\_64 (ports en progreso) & Redox Book; repos oficiales \\

Microsoft Windows & x86\_64, ARM64 & Microsoft Learn; Windows Internals \\

SerenityOS & x86\_64 (AArch64 en progreso) & Wiki/Repo oficial; guías build \\

FreeRTOS & MCU: ARM Cortex-M, RISC-V, AVR, ESP32, etc. & freertos.org; getting started; API \\

macOS (Darwin/XNU) & ARM64 (Apple Silicon); x86\_64 legado & Apple Open Source; docs dev \\

Haiku OS & x86\_64 (x86; ports ARM en trabajo) & Haiku User/Dev Guides; repos \\

Android (AOSP) & ARM/ARM64 (principal), x86/x86\_64 & source.android.com; dev docs \\
\hline
\end{muntab}

\end{landscape}

\begin{landscape}
\begin{muntab}{ll}{comparacionSO3b}{Comparación técnica de sistemas operativos - Comunidad}
\hline
\textbf{Sistema} & \textbf{Comunidad} \\
\hline
GNU/Linux & Muy grande, global \\

FreeBSD & Grande/estable \\

Minix 3 & Pequeña, enfocada en docencia \\

Redox OS & Pequeña/activa (investigación) \\

Microsoft Windows & Muy grande (ecosistema comercial) \\

SerenityOS & Pequeña/entusiasta, muy activa \\

FreeRTOS & Muy grande en embebidos \\

macOS (Darwin/XNU) & Grande (cerrado con OSS parcial) \\

Haiku OS & Mediana/pequeña; organizada \\

Android (AOSP) & Enorme (fabricantes/devs) \\
\hline
\end{muntab}

\end{landscape}